{
\section*{Resumen}\label{resumen}
\addcontentsline{toc}{chapter}{Resumen}
\thispagestyle{empty}
\let\clearpage\relax
Numerosos tipos de métodos de limpieza están disponibles en el mercado. Cada uno de estos tiene ventajas y desventajas, la energía solar es una alternativa viable para suplir el aumento de la demanda energética actual, ya que se tiene una disponibilidad del recurso que es inagotable y cuando produce energía no genera grandes impactos ambientales como se observa en recursos fósiles. La eficiencia y potencia entregada por un panel fotovoltaico son condiciones que se ven afectadas, principalmente por aspectos ambientales y tecnologías utilizadas en su proceso de fabricación, a esta problemática un limpiador es la solucion mas comun. Estos métodos son amplios y de diferentes métodos.\\
Apoyándonos del conocimiento cribado de la investigación previa para el proyecto se decidirá el método más eficiente y productivo para la limpieza e los paneles solares, una vez decidido se procederá a diseñar y construir el mecanismo que limpiara los paneles. Se adquirirá material que después pasará por diferentes procesos de manufactura para ser ensamblado en su totalidad con componentes de tipo mecánico, electrónico y de control. Para hacer pruebas de campo y confirmar su funcionamiento correcto.


\subsection*{palabras clave}\label{palabras}
\addcontentsline{toc}{chapter}{Palabras clave}
Limpiador, Panel Solar, Control, Automatización, Energía Limpia, Maculación.
\newpage
\section*{Abstract}\label{abstract}
\addcontentsline{toc}{chapter}{Abstract}
Numerous types of cleaning methods are available on the market. Each one of these has advantages and disadvantages, the solar energy is a viable alternative to the actual rising demand in energy,  we have a unlimited source of this resource that can produce energy without making an environmental impact like the fossil resources. The efficiency and power given by a photovoltaic panel are affected factors, mainly by environmental aspects  and  technologies used in their fabrication process, to this problematic the most common solution is a cleaner. this methods are large and different. \\
Relying on the cull knowledge of the previous research for the project, the most efficient and productive method for cleaning solar panels will be determined, once decided we will proceed to design and build the mechanism that will clean the panels. Material will be acquired that will then go through different manufacturing processes to be assembled entirely with mechanical, electronic and control components. To make field tests and confirm its correct operation.

\subsection*{keywords}\label{keywords}
\addcontentsline{toc}{chapter}{Keywords}
Cleaner, Solar Panel, Control, Automation, Clean Energy, Soiling.
}